\chapter{ادبیات تحقیق}

\section{مقدمه}

فصل دوم مفاهیم و معیارهای لازم برای تحلیل احساسات ضمنی را تعریف کرد. در این فصل، مسیر تحول پژوهش‌ها و نسبت آن‌ها با مسئلۀ پایان‌نامه بررسی می‌شود. هدف، تکرار تعریف‌ها نیست؛ بلکه باید روشن شود هر جریان پژوهشی کدام بخش مسئله را حل کرده و چه محدودیتی باقی گذاشته است.

مرور از روش‌های عمومی تحلیل احساسات آغاز می‌شود. سپس پژوهش‌های مبتنی بر جنبه و مجموعه‌داده‌های احساس ضمنی بررسی می‌شوند. در ادامه، مهندسی پرامپت، زنجیرۀ تفکر، خودسازگاری و دو روش مهم \lr{THOR} و \lr{SAoT} تحلیل می‌شوند. پایان فصل نیز به مقایسۀ روش‌ها و استخراج شکاف پژوهشی اختصاص دارد.

مقایسۀ عددی نتایج مقاله‌ها باید با احتیاط انجام شود. پژوهش‌ها از مدل‌های پایه، نسخۀ داده، شیوۀ نمونه‌گیری و تنظیمات متفاوت استفاده کرده‌اند. بنابراین، اعداد گزارش‌شده در یک مقاله لزوماً با اعداد مقالۀ دیگر قابل مقایسۀ مستقیم نیستند. تمرکز این فصل بر ایدۀ روش، نوع اطلاعات مصرف‌شده و سازوکار تولید تصمیم است.

\section{تحول رویکردهای تحلیل احساسات}

\subsection{روش‌های مبتنی بر واژگان و یادگیری ماشین}

پژوهش‌های نخستین تحلیل احساسات بیشتر بر واژه‌های دارای بار عاطفی تکیه داشتند. در این روش‌ها، متن با استفاده از فرهنگ واژگان، بسامد واژه‌ها یا ویژگی‌های سطحی نمایش داده می‌شد. سپس یک قاعده یا طبقه‌بند آماری قطبیت را تعیین می‌کرد \cite{pang2008opinion,wankhade2022survey}.

این رویکردها ساده، سریع و تا حدی تفسیرپذیر بودند. بااین‌حال، معنای واژه به بافت و هدف وابسته است. منفی‌سازها، تضاد، کنایه و چندمعنایی نیز رابطۀ ثابت میان واژه و قطبیت را مختل می‌کنند. این مشکل در تحلیل مبتنی بر جنبه جدی‌تر است، زیرا یک جمله ممکن است دربارۀ دو هدف، دو نگرش متفاوت بیان کند.

یادگیری ماشین نظارت‌شده بخشی از وابستگی به قواعد دستی را کاهش داد. ویژگی‌هایی مانند چندواژه، نقش دستوری و فاصلۀ واژه تا هدف به طبقه‌بند داده می‌شدند. بااین‌حال، کیفیت مدل همچنان به مهندسی ویژگی و شباهت میان داده‌های آموزش و آزمون وابسته بود. چنین بازنمایی‌هایی معمولاً رابطۀ دوربرد میان هدف و قرینۀ احساسی را ناقص ثبت می‌کردند.

این محدودیت در احساس ضمنی بنیادی‌تر است. ممکن است هیچ واژۀ مثبت یا منفی آشکاری در جمله وجود نداشته باشد. در این حالت، افزایش اندازۀ فرهنگ واژگان به تنهایی مسئله را حل نمی‌کند. مدل باید پیامد رویداد، انتظار متعارف یا رابطۀ آن با هدف را نیز استنتاج کند.

\subsection{روش‌های مبتنی بر شبکه‌های عصبی و مدل‌های زبانی}

شبکه‌های عصبی امکان یادگیری بازنمایی متن را بدون تعریف دستی همۀ ویژگی‌ها فراهم کردند. مدل‌های بازگشتی، پیچشی و سازوکار توجه برای تمرکز بر بخش‌های مرتبط با هدف به کار رفتند. در تحلیل مبتنی بر جنبه، نمایش هدف معمولاً همراه با بازنمایی جمله ساخته می‌شد تا قطبیت نسبت به همان هدف تعیین شود.

مدل‌های گرافی گام دیگری در این مسیر بودند. برای نمونه، شبکۀ توجه گراف رابطه‌ای\LTRfootnote{Relational Graph Attention Network} از رابطۀ نحوی میان واژه‌ها برای پیوند هدف و نشانه‌های نظر استفاده می‌کند \cite{wang2020rgat}. این روش‌ها وابستگی‌های ساختاری را بهتر از ویژگی‌های سطحی نمایش می‌دهند. بااین‌حال، نمودار نحوی لزوماً دانشی را که در متن بیان نشده است فراهم نمی‌کند.

مدل‌های زبانی ازپیش‌آموزش‌دیده، دانش زبانی گسترده‌تری را وارد مسئله کردند. سازگارسازی این مدل‌ها برای داده‌های برچسب‌خورده، عملکرد تحلیل مبتنی بر جنبه را بهبود داد. پژوهش‌های بعدی نیز اهداف پیش‌آموزش ویژه‌ای برای حساس‌کردن بازنمایی به هدف و قطبیت طراحی کردند \cite{li2021scapt}.

مدل‌های زبانی بزرگ، افزون بر بازنمایی، قابلیت تولید پاسخ و توضیح را در اختیار می‌گذارند. مسئله دیگر فقط نگاشت یک بردار به برچسب نیست. پژوهشگر می‌تواند صورت مسئله، فضای برچسب و مراحل تحلیل را با زبان طبیعی بیان کند. این تحول، مسیر استفاده از پرامپت و استدلال زنجیره‌ای را در تحلیل احساسات گشود.

\section{پژوهش‌های تحلیل احساسات مبتنی بر جنبه و ضمنی}

\subsection{مجموعه‌داده‌ها و چارچوب‌های ارزیابی}

رقابت \lr{SemEval 2014} یکی از مبناهای مهم تحلیل احساسات مبتنی بر جنبه است \cite{pontiki2014semeval}. این رقابت داده‌هایی از دامنه‌های لپ‌تاپ و رستوران ارائه کرد. وظایف آن شامل شناسایی جنبه و تعیین قطبیت نسبت به هدف بود. ساختار هدف‌محور این داده‌ها امکان بررسی جمله‌هایی را فراهم کرد که چند موضوع یا چند نگرش دارند.

بااین‌حال، داده‌های اصلی صرفاً برای مطالعۀ احساس ضمنی طراحی نشده بودند. لی و همکاران نمونه‌های صریح و ضمنی را از یکدیگر تفکیک کردند. آن‌ها همچنین پیش‌آموزش تقابلی نظارت‌شده\LTRfootnote{Supervised Contrastive Pre-training} را برای یادگیری احساس ضمنی در تحلیل مبتنی بر جنبه پیشنهاد کردند \cite{li2021scapt}. روش پیشنهادی آن‌ها با نام \lr{SCAPT} شناخته می‌شود. نسخۀ تفکیک‌شدۀ داده نیز در پژوهش‌های بعدی بازاستفاده شده است.

اهمیت این تفکیک در نوع شواهد موجود است. نمونه‌های صریح معمولاً نشانه‌ای نزدیک و آشکار برای نظر دارند. در مقابل، نمونۀ ضمنی مدل را وادار می‌کند از توصیف رویداد یا ویژگی به نگرش برسد. به همین دلیل، عملکرد مناسب روی کل دادۀ مبتنی بر جنبه الزاماً به معنای توانایی مناسب روی زیرمجموعۀ ضمنی نیست.

دو دامنۀ لپ‌تاپ و رستوران نیز الگوهای استنتاجی متفاوتی دارند. در رستوران، سرعت خدمت، زمان انتظار و کیفیت غذا با انتظارات روزمره پیوند دارند. در لپ‌تاپ، عمر باتری، وزن، قیمت و سازگاری نرم‌افزاری اهمیت بیشتری دارند. یک روش موفق باید هم ساختار مشترک مسئله و هم تفاوت دانش دامنه‌ای را در نظر بگیرد.

\subsection{روش‌های پیش از رویکردهای مولد}

روش \lr{SCAPT} با ساخت جفت‌های تقابلی، مدل را به تمایز نمونه‌های دارای قطبیت متفاوت هدایت می‌کند \cite{li2021scapt}. این جهت‌گیری از پیش‌آموزش عمومی هدفمندتر است. مدل می‌آموزد بازنمایی‌هایی بسازد که برای تشخیص احساس نسبت به جنبه مفید باشند. مزیت این روش آن است که دانش مسئله در پارامترها و فضای بازنمایی تثبیت می‌شود.

در کنار آن، روش‌های مبتنی بر ساختار نحوی تلاش کردند مسیر میان هدف و عبارت نظر را مدل کنند \cite{wang2020rgat}. مداخلۀ علّی نیز برای کاهش اتکای مدل به هم‌بستگی‌های سطحی پیشنهاد شد \cite{wang2022causal}. این پژوهش نشان می‌دهد بخشی از خطاهای احساس ضمنی از میان‌برهای آماری ناشی می‌شوند. مدل ممکن است به یک واژه یا الگوی پرتکرار تکیه کند، بدون آنکه رابطۀ واقعی شواهد و قطبیت را دریابد.

این خانواده از روش‌ها معمولاً به دادۀ برچسب‌خورده و آموزش ویژه نیاز دارد. خروجی آن‌ها نیز اغلب یک برچسب نهایی است. در نتیجه، تشخیص اینکه مدل کدام جنبه یا پیامد ضمنی را مبنای تصمیم قرار داده دشوار می‌شود. همچنین، انتقال به دامنه‌ای تازه می‌تواند به سازگارسازی مجدد نیاز داشته باشد.

ورود مدل‌های مولد، پرسش پژوهش را تغییر داد. اکنون می‌توان از مدل خواست رابطۀ پنهان میان هدف، ویژگی ضمنی و نگرش را به زبان طبیعی آشکار کند. این قابلیت جای روش‌های بازنمایی‌محور را به طور کامل نمی‌گیرد، اما ابزار تازه‌ای برای صورت‌بندی و بررسی فرایند استنتاج فراهم می‌کند.

\subsection{چالش‌های مشترک گزارش‌شده}

نخستین چالش، نبود قرینۀ احساسی صریح است. مدل باید بفهمد یک رویداد برای هدف مطلوب، نامطلوب یا بی‌اثر است. این داوری به دانش عمومی و انتظارهای دامنه‌ای وابسته است. یک رویداد واحد نیز ممکن است برای دو هدف پیامد متفاوتی داشته باشد.

چالش دوم، نسبت‌دادن احساس به هدف درست است. جمله ممکن است چند موجودیت یا چند ویژگی داشته باشد. اگر مدل شواهد مربوط به یک هدف را به هدف دیگر منتقل کند، قطبیت نهایی نادرست می‌شود. روش‌های نحوی این مسئله را کاهش می‌دهند، اما استنتاج ضمنی همیشه با رابطۀ نحوی مستقیم هم‌راستا نیست.

چالش سوم، تمایز خنثی از احساس ضعیف و ضمنی است. نبود واژۀ ارزیابانه می‌تواند مدل را به پیش‌بینی خنثی سوق دهد. از سوی دیگر، جست‌وجوی افراطی برای معنای پنهان ممکن است جملۀ واقعاً خنثی را احساسی جلوه دهد. بنابراین، روش باید هم از حذف شواهد ضمنی و هم از تفسیر بیش‌ازحد جلوگیری کند.

چالش چهارم، نبود تضمین برای وفاداری توضیح است. یک مدل مولد می‌تواند پس از انتخاب برچسب، توضیحی ظاهراً قانع‌کننده بسازد. پس وجود متن استدلال به تنهایی اثبات نمی‌کند که تصمیم از همان مسیر حاصل شده است. ارزیابی برچسب، سازگاری مراحل و منشأ تصمیم باید از یکدیگر تفکیک شوند.

\section{مهندسی پرامپت}

مرور روش‌های پرامپت نشان می‌دهد که صورت‌بندی وظیفه، نمونه‌های راهنما و قالب پاسخ بر رفتار مدل اثر می‌گذارند \cite{liu2023prompting}. در تحلیل احساسات، یک پرامپت حداقلی فقط جمله، هدف و برچسب‌های مجاز را مشخص می‌کند. پرامپت غنی‌تر می‌تواند تعریف قطبیت، قید هدف‌محوری یا مراحل تحلیل را نیز اضافه کند.

پرامپت مستقیم یک خط پایۀ ضروری است. این روش نشان می‌دهد مدل بدون الزام به تولید توضیح تا چه اندازه مسئله را حل می‌کند. اگر روش استدلالی با خط پایۀ مستقیم مقایسه نشود، نمی‌توان اثر واقعی مراحل افزوده را سنجید. همچنین، هر افزایش عملکرد را نمی‌توان صرفاً به طولانی‌ترشدن پرامپت نسبت داد.

قالب خروجی نیز اهمیت دارد. پاسخ آزاد ممکن است حاوی چند برچسب، توضیح مبهم یا عبارتی خارج از فضای برچسب باشد. قالب ساختاریافته، استخراج خودکار و تشخیص پاسخ نامعتبر را آسان‌تر می‌کند. بااین‌حال، سخت‌گیری قالب نباید به اندازه‌ای باشد که ظرفیت تحلیل مدل را محدود کند.

پرامپت صفرنمونه هزینه آماده‌سازی مثال را کاهش می‌دهد، اما به تفسیر مدل از دستور وابسته است. پرامپت چندنمونه می‌تواند مرز میان برچسب‌ها را روشن کند. بااین‌حال، گزینش مثال‌ها ممکن است سوگیری ایجاد کند. در هر دو حالت، متن دستور، ترتیب اطلاعات و نحوۀ نمایش هدف باید ثابت و مستند باشند.

در احساس ضمنی، پرامپت خوب باید مدل را به شواهد مرتبط با هدف هدایت کند. دستور کلی «احساس را تشخیص بده» ممکن است نگرش کلی جمله را جایگزین نگرش هدف‌محور کند. ازاین‌رو، پژوهش‌های جدیدتر وظیفه را به پرسش‌های کوچک‌تر تقسیم کرده‌اند.

\section{زنجیرۀ تفکر در پژوهش‌های پردازش زبان طبیعی}

وی و همکاران نشان دادند که افزودن نمونه‌های دارای مراحل میانی می‌تواند توانایی استدلال مدل‌های زبانی بزرگ را آشکار کند \cite{wei2022cot}. ایدۀ اصلی آن‌ها تولید چند گام زبان طبیعی پیش از پاسخ نهایی بود. این روش ابتدا در مسئله‌های حسابی، نمادین و استدلال عقل سلیم بررسی شد.

اهمیت این پژوهش برای تحلیل احساسات ضمنی در شباهت ساختاری وظایف است. در هر دو حالت، پاسخ نهایی از مشاهدۀ یک نشانۀ صریح به دست نمی‌آید. مدل باید اطلاعات متن را با دانش زمینه‌ای ترکیب و رابطۀ چندمرحله‌ای را دنبال کند. بااین‌حال، مراحل مناسب برای هر وظیفه باید بر اساس ساختار همان مسئله طراحی شوند.

زنجیرۀ تفکر را نباید با صحت تضمین‌شده یکی دانست. مراحل تولیدی می‌توانند حاوی فرض نادرست، جابه‌جایی هدف یا تناقض باشند. طولانی‌ترشدن استدلال نیز الزاماً کیفیت را افزایش نمی‌دهد. هر مرحلۀ اضافی یک فرصت تازه برای اصلاح، اما هم‌زمان یک نقطۀ تازه برای خطا ایجاد می‌کند.

در ادبیات تحلیل احساسات ضمنی، این ایده از یک دستور عمومی به تجزیۀ وظیفه تبدیل شد. به جای درخواست مبهم برای «فکرکردن مرحله‌به‌مرحله»، پژوهش‌ها پرسش‌هایی دربارۀ جنبه، نظر و رابطۀ آن‌ها با قطبیت مطرح کردند. روش \lr{THOR} نمونۀ شاخص این صورت‌بندی ساختاریافته است.

\section{روش \lr{THOR} برای تحلیل احساسات ضمنی}

فی و همکاران روش \lr{THOR} را برای استدلال دربارۀ احساس ضمنی ارائه کردند \cite{fei2023thor}. این روش مسئله را به سه پرسش وابسته تقسیم می‌کند. هر پاسخ، زمینۀ لازم برای مرحلۀ بعد را فراهم می‌سازد. هدف آن است که مدل پیش از انتخاب قطبیت، اجزای پنهان رابطۀ احساسی را آشکار کند.

\subsection{تجزیۀ استدلال به جنبه، نظر و قطبیت}

مرحلۀ نخست، جنبه یا ویژگی ضمنی مرتبط با هدف را استنتاج می‌کند. ممکن است جمله نام هدف را داشته باشد، اما ویژگی ارزیابی‌شده را آشکار بیان نکند. تشخیص این ویژگی فضای تفسیر را محدود می‌کند. برای مثال، توصیف مدت کارکرد دستگاه می‌تواند به جنبه‌ای مانند باتری مربوط باشد.

مرحلۀ دوم، نظر ضمنی نسبت به جنبه را استخراج می‌کند. این نظر لزوماً یک صفت موجود در جمله نیست. مدل باید پیامد ویژگی یا رویداد را با انتظارات رایج مقایسه کند. پاسخ مرحلۀ دوم به متن، هدف و جنبه‌ای که در مرحلۀ نخست ساخته شده وابسته است.

مرحلۀ سوم، قطبیت را از رابطۀ میان هدف، جنبه و نظر تعیین می‌کند. این ترتیب، جهش مستقیم از جمله به برچسب را به چند تصمیم قابل مشاهده تبدیل می‌کند. در نتیجه، پژوهشگر می‌تواند محل احتمالی انحراف را دقیق‌تر بررسی کند.

وابستگی مراحل هم مزیت و هم محدودیت است. اگر جنبه درست تشخیص داده شود، مرحلۀ نظر راهنمای متمرکزتری دریافت می‌کند. اما خطای مرحلۀ نخست می‌تواند به مراحل بعد منتقل شود. بنابراین، صحت هر گام و سازگاری میان گام‌ها اهمیت دارد.

\subsection{خودسازگاری و بازبینی در \lr{THOR}}

روش \lr{THOR} برای کاهش ناپایداری تولید از نمونه‌گیری چندباره و رأی‌گیری استفاده می‌کند. ایدۀ خودسازگاری آن است که مسیرهای استدلال متنوع تولید شوند و پاسخ پرتکرار انتخاب شود \cite{wang2023selfconsistency}. این سازوکار وابستگی نتیجه به یک اجرای تصادفی را کاهش می‌دهد.

در صورت وجود دادۀ نظارت‌شده، مقالۀ \lr{THOR} سازوکاری برای بازبینی استدلال نیز مطرح می‌کند. نمونۀ خطادار می‌تواند همراه با مسیر اصلاح‌شده به مدل ارائه شود. این بخش نشان می‌دهد که تولید استدلال و اصلاح استدلال دو مسئلۀ جدا هستند. بااین‌حال، دسترسی به نمونۀ اصلاح‌شده یا نظارت مناسب همیشه ممکن نیست.

خودسازگاری نیز تضمین‌کنندۀ پاسخ درست نیست. اگر یک سوگیری در بیشتر مسیرها تکرار شود، رأی اکثریت همان خطا را تثبیت می‌کند. افزون بر این، چند اجرای یک روش ممکن است تنوع معنایی کافی نداشته باشند. بنابراین، توافق باید به عنوان نشانۀ ثبات تفسیر شود، نه اثبات صحت.

\subsection{نقاط قوت و محدودیت‌های \lr{THOR}}

نقطۀ قوت اصلی \lr{THOR}، انطباق مراحل استدلال با ساختار تحلیل احساسات مبتنی بر جنبه است. روش به طور مشخص می‌پرسد چه چیزی درباره هدف ارزیابی شده و نظر پنهان چیست. این طراحی از یک دستور عمومی دقیق‌تر است و ردپایی برای تحلیل خطا فراهم می‌کند.

نقطۀ قوت دوم، امکان تفکیک خطاها است. خطا می‌تواند از تشخیص جنبه، استنتاج نظر یا نگاشت نظر به قطبیت ناشی شود. این تفکیک برای طراحی بازبین و تحلیل کیفی مفید است. همچنین، پاسخ استدلالی می‌تواند در کنار پیش‌بینی مستقیم به عنوان منبعی مکمل نگهداری شود.

محدودیت نخست، انتشار خطا در زنجیره است. هر مرحلۀ بعدی به پاسخ‌های پیشین وابسته است. محدودیت دوم، هزینۀ بیشتر اجرای چندمرحله‌ای و چندباره است. محدودیت سوم، امکان تولید توضیحی روان اما ناسازگار با شواهد جمله است.

افزون بر این، مقالۀ اصلی عمدتاً کیفیت یک مسیر استدلالی را ارزیابی می‌کند. اختلاف میان پاسخ مستقیم و پاسخ استدلالی خود به مسئله‌ای مستقل تبدیل می‌شود. اگر یکی از دو منبع در گروهی از نمونه‌ها بهتر باشد، انتخاب ثابت یک منبع می‌تواند بخشی از اطلاعات مفید را از بین ببرد.

\section{بازبینی و خوداصلاحی در تحلیل احساسات ضمنی}

بازبینی از مدل می‌خواهد پاسخ اولیه را دوباره بررسی کند. این کار می‌تواند با پرسشی کلی یا مجموعه‌ای از معیارهای مشخص انجام شود. بازبینی کلی انعطاف‌پذیر است، اما خروجی آن ممکن است مبهم باشد. بازبینی ساختاریافته، نوع ناسازگاری و پیشنهاد اصلاح را جدا می‌کند.

دوان و وانگ روش \lr{SAoT} را برای تحلیل احساسات ضمنی پیشنهاد کردند \cite{duan2024isa}. این روش ابتدا جنبه‌ها و نظرهای ضمنی را تحلیل می‌کند. سپس مدل تحلیل خود را بازنگری می‌کند و در پایان قطبیت را نتیجه می‌گیرد. در این طراحی، بازاندیشی بخشی صریح از زنجیرۀ تصمیم است.

تفاوت مهم \lr{SAoT} و \lr{THOR} در جایگاه بازبینی است. \lr{THOR} بر تجزیۀ گام‌به‌گام و خودسازگاری تأکید دارد. \lr{SAoT} یک مرحلۀ بازاندیشی را پیش از نتیجه‌گیری نهایی وارد می‌کند. هر دو روش می‌کوشند استنتاج پنهان را آشکار کنند، اما کنترل خطا را به شکل متفاوتی صورت‌بندی می‌کنند.

بازبینی می‌تواند خطای اولیه را اصلاح کند، ولی ممکن است پاسخ درست را نیز تغییر دهد. مدل بازبین همان محدودیت‌های دانش و سوگیری مدل تولیدکننده را دارد. اگر دستور بازبینی صرفاً مدل را به تغییر پاسخ ترغیب کند، اصلاح ظاهری جای بررسی شواهد را می‌گیرد. بنابراین، فعال‌سازی بازبینی و پذیرش پیشنهاد آن باید از هم جدا باشند.

پژوهش‌های یادشده نشان می‌دهند که «دوباره فکرکردن» به تنهایی معیار تصمیم نیست. خروجی بازبینی باید قابل استخراج و قابل مقایسه با منابع دیگر باشد. همچنین، لازم است موارد توافق، اختلاف و نبود شواهد کافی از یکدیگر تفکیک شوند.

\section{خودسازگاری در استدلال چندمرحله‌ای}

وانگ و همکاران خودسازگاری را برای بهبود استدلال زنجیره‌ای معرفی کردند \cite{wang2023selfconsistency}. به جای انتخاب نخستین مسیر، چند مسیر با نمونه‌گیری تصادفی تولید می‌شوند. سپس پاسخ نهایی با تجمیع پاسخ‌ها تعیین می‌شود. این روش از تنوع مسیرهای ممکن برای کاهش خطای یک اجرای منفرد استفاده می‌کند.

خودسازگاری یک روش استدلال مستقل نیست. کیفیت آن به پرامپت، مدل پایه و تنوع نمونه‌گیری وابسته است. اگر همۀ اجراها از یک برداشت نادرست آغاز شوند، رأی‌گیری کمکی نمی‌کند. اگر پاسخ‌ها بسیار پراکنده باشند، رأی اکثریت نیز ممکن است شکننده باشد.

در تحلیل احساسات ضمنی، دو سطح تجمیع قابل تصور است. در سطح نخست، برچسب نهایی هر اجرای کامل رأی می‌دهد. در سطح دوم، پاسخ هر مرحلۀ استدلال جداگانه تجمیع می‌شود. سطح دوم می‌تواند خطای میانی را زودتر مهار کند، اما به نرمال‌سازی پاسخ‌های زبانی نیاز دارد.

رأی‌گیری همچنین اطلاعاتی فراتر از برچسب فراهم می‌کند. نسبت آرا می‌تواند نشانۀ ثبات باشد و مسیرهای مخالف می‌توانند برای تحلیل اختلاف نگهداری شوند. بااین‌حال، این نشانه نباید بدون کالیبراسیون به احتمال صحت تعبیر شود.

\section{کنترل‌گرها و انتخاب منبع پیش‌بینی}

بیشتر روش‌های مرورشده یک مسیر اصلی را برای تولید پاسخ نهایی تعیین می‌کنند. روش بازنمایی‌محور از طبقه‌بند، روش مستقیم از پاسخ کوتاه و روش زنجیره‌ای از نتیجۀ استدلال استفاده می‌کند. با وجود این، هیچ مسیر واحدی در همۀ نمونه‌ها برتری تضمین‌شده ندارد.

پیش‌بینی مستقیم معمولاً کوتاه‌تر و کم‌هزینه‌تر است. این مسیر گاهی از دانش درونی مدل به شکل مؤثری استفاده می‌کند، بدون آنکه در مراحل میانی منحرف شود. مسیر استدلالی در نمونه‌هایی مفید است که رابطۀ رویداد و قطبیت به توضیح چندمرحله‌ای نیاز دارد. بازبینی نیز ممکن است ناسازگاری مشخصی را آشکار کند.

این تفاوت، مسئلۀ انتخاب منبع را مطرح می‌کند. کنترل‌گر باید شواهدی مانند توافق، اعتبار قالب و نشانه‌های خطا را بررسی کند. سپس یک سیاست انتخاب تعیین می‌کند کدام منبع پاسخ نهایی را تأمین کند. این نگاه با انتخاب همیشگی پاسخ طولانی‌تر یا آخرین پاسخ متفاوت است.

ادبیات \lr{THOR} و \lr{SAoT} اهمیت مراحل میانی و بازبینی را نشان می‌دهد \cite{fei2023thor,duan2024isa}. بااین‌حال، این پژوهش‌ها انتخاب داده‌محور میان پیش‌بینی مستقیم، استدلالی و تشخیصی را به عنوان مسئلۀ محوری صورت‌بندی نمی‌کنند. ازاین‌رو، مدیریت اختلاف منابع همچنان جای بررسی دارد.

سیاست انتخاب باید بدون استفاده از برچسب طلایی آزمون ساخته شود. در غیر این صورت، انتخاب بهترین منبع پس از مشاهدۀ پاسخ درست به برآوردی غیرواقعی منجر می‌شود. همچنین، رفتار سیاست در وضعیت‌های دیده‌نشده باید روشن باشد. جزئیات سیاست این پایان‌نامه در فصل چهارم ارائه می‌شود.

\section{مقایسۀ پژوهش‌های مرتبط}

جدول \ref{tab:related-work-comparison} پژوهش‌های محوری را از نظر هدف، سازوکار و محدودیت مرتبط مقایسه می‌کند. منظور از محدودیت، نفی دستاورد پژوهش نیست. این ستون نشان می‌دهد برای ساخت یک سامانۀ چندمنبعی چه مسئله‌ای همچنان باقی می‌ماند.

\begin{table}[htbp]
    \centering
    \caption{مقایسۀ پژوهش‌های محوری مرتبط با تحلیل احساسات ضمنی.}
    \label{tab:related-work-comparison}
    \small
    \setlength{\tabcolsep}{5pt}
    \renewcommand{\arraystretch}{1.25}
    \begin{tabular}{|p{3.0cm}|p{5.2cm}|p{5.2cm}|}
        \hline
        {\raggedleft\textbf{پژوهش}\par} &
        {\raggedleft\textbf{ایده و سازوکار}\par} &
        {\raggedleft\textbf{محدودیت مرتبط}\par}
        \tabularnewline
        \hline
        {\raggedleft لی و همکاران؛ \lr{SCAPT} \cite{li2021scapt}\par} &
        {\raggedleft پیش‌آموزش تقابلی نظارت‌شده برای یادگیری بازنمایی حساس به احساس ضمنی\par} &
        {\raggedleft نیازمند آموزش ویژه و فاقد ردپای زبانی برای استنتاج\par}
        \tabularnewline
        \hline
        {\raggedleft وانگ و همکاران \cite{wang2022causal}\par} &
        {\raggedleft مداخلۀ علّی برای کاهش اتکا به هم‌بستگی‌های سطحی\par} &
        {\raggedleft تمرکز بر بازنمایی یک مدل، نه انتخاب میان منابع\par}
        \tabularnewline
        \hline
        {\raggedleft وی و همکاران؛ زنجیرۀ تفکر \cite{wei2022cot}\par} &
        {\raggedleft تولید مراحل زبانی میانی پیش از پاسخ نهایی\par} &
        {\raggedleft مراحل عمومی و غیرتخصصی برای احساس ضمنی\par}
        \tabularnewline
        \hline
        {\raggedleft وانگ و همکاران؛ خودسازگاری \cite{wang2023selfconsistency}\par} &
        {\raggedleft نمونه‌گیری چند مسیر و انتخاب پاسخ پرتکرار\par} &
        {\raggedleft امکان تثبیت خطای مشترک با رأی اکثریت\par}
        \tabularnewline
        \hline
        {\raggedleft فی و همکاران؛ \lr{THOR} \cite{fei2023thor}\par} &
        {\raggedleft استنتاج مرحله‌ای جنبه، نظر و قطبیت همراه با خودسازگاری\par} &
        {\raggedleft انتشار خطا میان مراحل و نبود انتخاب تطبیقی منبع\par}
        \tabularnewline
        \hline
        {\raggedleft دوان و وانگ؛ \lr{SAoT} \cite{duan2024isa}\par} &
        {\raggedleft تحلیل جنبه و نظر، بازاندیشی و سپس استنتاج قطبیت\par} &
        {\raggedleft امکان تغییر پاسخ درست و نبود سیاست مستقل برای پذیرش اصلاح\par}
        \tabularnewline
        \hline
    \end{tabular}
\end{table}

روند کلی جدول از یادگیری بازنمایی به آشکارسازی استدلال حرکت می‌کند. \lr{SCAPT} و روش علّی بر بهبود مدل آموزش‌دیده تمرکز دارند. زنجیرۀ تفکر و \lr{THOR} ساخت تصمیم را به زبان طبیعی قابل مشاهده می‌کنند. خودسازگاری و \lr{SAoT} نیز به ترتیب ثبات و بازبینی را وارد فرایند می‌کنند.

بااین‌حال، هیچ‌یک از این مؤلفه‌ها به تنهایی همۀ نیازها را پوشش نمی‌دهد. بازنمایی قوی ممکن است توضیحی نداشته باشد. توضیح چندمرحله‌ای ممکن است منحرف شود. رأی اکثریت ممکن است خطای مشترک را تقویت کند و بازبینی ممکن است پاسخ درست را تغییر دهد. این ناهمگونی، انگیزۀ ترکیب کنترل‌شدۀ منابع را فراهم می‌کند.

\section{شکاف تحقیقاتی و جایگاه پژوهش حاضر}

مرور ادبیات چهار شکاف مرتبط را نشان می‌دهد. نخست، بسیاری از روش‌ها یک مسیر را پاسخ نهایی و قطعی فرض می‌کنند. درحالی‌که پیش‌بینی مستقیم و استدلالی می‌توانند در نمونه‌های متفاوت برتری داشته باشند. اختلاف آن‌ها صرفاً مزاحمت نیست و می‌تواند اطلاعات تشخیصی فراهم کند.

دوم، بازبینی اغلب به شکل یک مرحلۀ تولیدی دیگر تعریف می‌شود. تولید پیشنهاد اصلاح با پذیرش آن یکسان نیست. سامانه باید تشخیص دهد بازبینی چه نوع خطایی را گزارش کرده و تا چه اندازه قابل اتکاست. بدون این جداسازی، خوداصلاحی می‌تواند به خودتخریبی تبدیل شود.

سوم، خودسازگاری عمدتاً مسیرهای یک روش را تجمیع می‌کند. این سازوکار برای کاهش نوسان مفید است، اما تفاوت میان منابع ناهمگون را مدل نمی‌کند. پاسخ مستقیم، استدلال زنجیره‌ای و بازبینی تشخیصی ماهیت یکسانی ندارند. رأی برابر میان آن‌ها نیز لزوماً بهترین سیاست نیست.

چهارم، انتخاب پس از مشاهدۀ برچسب آزمون فاقد اعتبار تجربی است. لازم است سیاست انتخاب فقط با داده‌های مجاز ساخته شود و برای وضعیت ناشناخته رفتار مشخصی داشته باشد. همچنین، منشأ هر پاسخ نهایی باید قابل ردیابی بماند تا سهم واقعی هر منبع ارزیابی شود.

پژوهش حاضر در همین نقطه قرار می‌گیرد. هدف، ابداع یک مدل زبانی یا روش بنیادی تازه برای زنجیرۀ تفکر نیست. این پایان‌نامه پیش‌بینی مستقیم، استدلال چندمرحله‌ای و بازبینی تشخیصی را منابع مکمل در نظر می‌گیرد. سپس مسئلۀ اصلی را به انتخاب کنترل‌شدۀ منبع نهایی تبدیل می‌کند.

این جایگاه با \lr{THOR} و \lr{SAoT} پیوند دارد، اما با آن‌ها یکسان نیست. از \lr{THOR} ساختار استدلال جنبه، نظر و قطبیت الهام گرفته می‌شود. ایدۀ بازبینی نیز برای تشخیص ناسازگاری به کار می‌رود. تفاوت اصلی در جداسازی تصمیم میانی از منبع نهایی و ارزیابی سیاست انتخاب است.

جزئیات این سیاست در این فصل ارائه نمی‌شود. فصل چهارم معماری، مسیرهای پیش‌بینی، بازبین و سازوکار انتخاب را به طور دقیق شرح می‌دهد. فصل پنجم نیز نشان می‌دهد هر مؤلفه و هر منبع در داده‌های آزمایشی چه رفتاری داشته است.

\section{جمع‌بندی}

این فصل سیر پژوهش را از ویژگی‌های سطحی به بازنمایی‌های آموخته‌شده و سپس استدلال مولد مرور کرد. روش‌های واژه‌محور و آماری برای نشانه‌های آشکار مناسب‌تر بودند. مدل‌های عصبی، گرافی و پیش‌آموزش تقابلی رابطۀ هدف و متن را بهتر مدل کردند، اما همچنان به آموزش ویژه وابسته بودند.

زنجیرۀ تفکر امکان آشکارسازی مراحل استنتاج را فراهم کرد. \lr{THOR} این ایده را با تجزیۀ جنبه، نظر و قطبیت برای احساس ضمنی تخصصی کرد. خودسازگاری نوسان اجرا را کاهش داد و \lr{SAoT} بازاندیشی را وارد زنجیره کرد. در عین حال، انتشار خطا، توضیح غیرهمسو و اصلاح نادرست همچنان باقی ماندند.

نتیجۀ اصلی مرور آن است که هیچ منبعی همواره برتر نیست. مسئلۀ پژوهش حاضر فقط تولید یک استدلال دیگر نیست؛ بلکه باید اختلاف منابع را مدیریت و منشأ پاسخ نهایی را روشن کند. فصل بعد بر پایۀ این شکاف، کار پیشنهادی و اجزای اجرایی سامانه را تشریح می‌کند.
